\documentclass[tikz,border=6pt]{standalone}
\usepackage{gensymb}
\usetikzlibrary{arrows.meta,calc,angles,patterns,quotes}

\begin{document}
\begin{tikzpicture}

\pgfdeclareimage{silhouetteAvion}{DR400.pdf}

% Paramètres (0° = Nord, sens horaire)
\newcommand{\vitesseVent}{50}
\newcommand{\vitesseMaxVent}{50}
\newcommand{\vitessePropreAvion}{120}
\newcommand{\directionVent}{65}
\newcommand{\routeSuivie}{90}

% Dérive et cap suivi
\pgfmathsetmacro{\facteurDeBase}{60/\vitessePropreAvion}
\pgfmathsetmacro{\deriveMax}{\facteurDeBase*\vitesseVent}
\pgfmathsetmacro{\derive}{90-\deriveMax*sin(\routeSuivie-\directionVent)}
\newcommand{\capSuivi}{\derive}

% Conversion caps/route/vent -> angles TikZ
\pgfmathsetmacro{\angleCapTikz}{90-\capSuivi}
\pgfmathsetmacro{\angleRouteTikz}{90-\routeSuivie}
\pgfmathsetmacro{\angleVentTikz}{-(\directionVent+90)} % flux (vient de directionVent)

% Longueurs
\pgfmathsetmacro{\longueurVent}{2*\vitesseVent/\vitesseMaxVent}
\pgfmathsetmacro{\ventFace}{\vitesseVent*cos(\directionVent-\capSuivi)}
\pgfmathsetmacro{\vitesseSol}{\vitessePropreAvion - \ventFace}
\pgfmathsetmacro{\longueurRoute}{8*\vitesseSol/\vitessePropreAvion} % x2
\pgfmathsetmacro{\longueurCap}{8}                                  % x2

% Points clés
\coordinate (centreFlecheVent)     at (-8.5,-1);
\coordinate (centreSilhouetteAvion) at (-8.5,-4);

% Coordonnées finales
\coordinate (debutFlecheVent) at ($(centreFlecheVent)+(\angleVentTikz-180:\longueurVent/2)$);
\coordinate (finFlecheVent)   at ($(centreFlecheVent)+(\angleVentTikz:\longueurVent/2)$);
\coordinate (finCap)          at ($(centreSilhouetteAvion)+(\angleCapTikz:\longueurCap)$);
\coordinate (finRoute)        at ($(centreSilhouetteAvion)+(\angleRouteTikz:\longueurRoute)$);

% Cadre
\draw (-10,-8) rectangle (4,0);

% Flèche de vent
\draw[-{Triangle[width=18pt,length=8pt]}, line width=10pt, blue]
  (debutFlecheVent) -- (finFlecheVent)
  node[midway, sloped, above]{\vitesseVent\ kts};

% Vecteur cap suivi (rouge, x2) – étiquette placée au bout, côté externe
\draw[->, red, thick, line width=4pt]
  (centreSilhouetteAvion) -- (finCap)
  node[above right=3pt and 1pt] {\vitessePropreAvion\ kts}
  node (vecteurCapSuivi) [black,right] {$\pgfmathprintnumber[fixed,precision=1]{\capSuivi}\degree$};

% Vecteur route suivie (noir) – étiquette au bout, côté externe
\draw[->, black, thick, line width=2pt]
  (centreSilhouetteAvion) -- (finRoute)
  node[below right=3pt and 1pt] {\pgfmathprintnumber[fixed,precision=1]{\vitesseSol}\ kts}
  node (vecteurRouteSuivie) [black,right] {$\routeSuivie\degree$};

% Silhouette avion
\node [scale=0.25, rotate=\routeSuivie-\capSuivi] at (centreSilhouetteAvion)
  {\pgfbox[center,center]{\pgfuseimage{silhouetteAvion}}};

% Arc de dérive (toujours le plus petit) – affichage à 1 décimale
\pgfmathsetmacro{\delta}{mod(\routeSuivie-\capSuivi+360,360)}
\ifdim \delta pt > 180pt
  \pgfmathsetmacro{\valeurAlpha}{360-\delta}
  \pic [draw,<->,"$\alpha = \pgfmathprintnumber[fixed,precision=1]{\valeurAlpha}\degree$", angle eccentricity=1.4, angle radius=2cm]
      {angle = vecteurCapSuivi--centreSilhouetteAvion--vecteurRouteSuivie};
\else
  \pgfmathsetmacro{\valeurAlpha}{\delta}
  \pic [draw,<->,"$\alpha = \pgfmathprintnumber[fixed,precision=1]{\valeurAlpha}\degree$", angle eccentricity=1.4, angle radius=2cm]
      {angle = vecteurRouteSuivie--centreSilhouetteAvion--vecteurCapSuivi};
\fi

\end{tikzpicture}
\end{document}